\documentclass[10pt,a4paper]{article}
\usepackage[margin=1.0cm, top=1.0cm, bottom=1.0cm, left=1.0cm, right=1.0cm]{geometry} % Margenes del documento
\usepackage{titlesec} % Formato de títulos o encabezados
\usepackage{hyperref}              % Para enlaces a LinkedIn o Email
\usepackage{xcolor}                % Para un toque de color sutil
\usepackage{custom}
\usepackage[utf8]{inputenc}
\usepackage[T1]{fontenc}    % Esencial para que las fuentes se rendericen bien
\usepackage[spanish]{babel}

% --- Configuración de párrafos --- %
\setlength{\parskip}{5pt}

% --- Configuración de colores --- %
\definecolor{accent}{RGB}{52, 152, 219}

% --- Configuración global de \section --- %
\titleformat{\section}
{\normalfont\large\bfseries\uppercase} % 1. MAYÚSCULAS y Negrita
{}                                     % 2. Quitar el bullet numérico (etiqueta vacía)
{0pt}                                  % 3. Sin espacio de numeración
{}									% 4. Linea horizontal SUPERIOR
[%
\vspace{-1ex}\rule{\textwidth}{0.4pt} % 5. Línea horizontal INFERIOR
]

% --- Configuración de hipervínculos --- %
\hypersetup{
	colorlinks=true,      % Activa el color en el texto del link
	linkcolor=accent,     % Color para enlaces internos
	filecolor=accent,     % Color para archivos
	urlcolor=accent,      % Color para URLs externas (LinkedIn, etc)
	citecolor=accent,
	pdfborder={0 0 0}     % Elimina definitivamente cualquier recuadro
}

\begin{document}
	
	\cvTitle{Ramiro  Martínez D'Elía}
	\cvSubtitle{Software Engineer - Solution Architect}
	
	\section{Acerca de mi}

	
	Software Engineer con +10 años de trayectoria, especializado en el diseño y modernización de arquitecturas orientadas a servicios y la transición de sistemas legacy hacia microservicios escalables. Mi enfoque integra la gestión estratégica del backlog con una sólida capacidad de ejecución en todo el ciclo de vida del software.
	
	Poseo experiencia en el diseño e implementación de soluciones de alta disponibilidad para procesamiento  masivo de datos, con especial énfasis en la resiliencia de sistemas mediante la observabilidad proactiva (utilizando herramientas del tipo APM).
	
	Tengo capacidad para traducir dominios de negocio complejos, como la liquidación impositiva y flujos financieros, en modelos técnicos robustos, priorizando la calidad del código y la gestión inteligente de la deuda técnica para asegurar la escalabilidad a largo plazo.
	
	\section{Contacto}
	
	\begin{center}
		% TODO: No me gusta la presentación
		\small
		\textbf{Teléfono:} +54 221 6393338 \vline\hfill
		\textbf{Email:} \href{mailto:ramiro.md90@gmail.com}{ramiro.md90@gmail.com} \vline\hfill
		\textbf{LinkedIn:} \href{https://linkedin.com/in/ramiromd}{/in/ramiromd} \vline\hfill
		\textbf{GitHub:} \href{https://github.com/ramiromd}{/ramiromd} \vline\hfill
		\textbf{Ciudad:} La Plata, Provincia de Buenos Aires, Argentina
	\end{center}
	
	\section{Habilidades}
	{
		\setlength{\parindent}{0pt}
		
		\textbf{Arquitectura y calidad:} Microservicios, Programación Orientada a Objetos (POO), Principios SOLID, Clean Architecture, Patrones de diseño, Event sourcing.
		
		\textbf{Lenguajes:} PHP, Java, JavaScript.
		
		\textbf{Frameworks:} Spring Boot, Symfony.
		
		\textbf{Bases de datos:} Relacionales (MySQL, PostgreSQL) y NoSQL (MongoDB, Redis).
		
		\textbf{DevOps y observabilidad:} Docker, Kubernetes, CI/CD, Elastic (APM).
		
		\textbf{Soft Skills:}  Visión estratégica de producto, Liderazgo técnico y mentoría, gestión de deuda técnica, gestión de trade-offs y deuda técnica, comunicación con stakeholders, traducción de complejidad normativa, pensamiento analítico.
	}
	
	\section{Experiencia profesional}
	
	\cvExperience{Solution Architect - Geopagos}{Marzo 2025 - Actualidad}{Ciudad Autónoma de Buenos Aires, Argentina} % TODO: Agregar equipo
	
	\noindent\textbf{Responsabilidades}
	\begin{itemize}
		\item Liderar la priorización del backlog técnico y de producto, equilibrando la entrega de funcionalidades con la resolución de deuda técnica y la escalabilidad de la plataforma.
		\item Definir el roadmap para la transición de flujos legacy hacia una arquitectura orientada a servicios, priorizando la resiliencia y el mantenimiento a largo plazo.
		\item Actuar como nexo técnico-funcional con socios externos y líderes de otros equipos para asegurar la consistencia de los servicios transversales.
		\item Coordinar la alineación de requerimientos complejos con líderes de producto y técnica de otros equipos, asegurando la consistencia funcional y técnica de los servicios transversales a nivel plataforma.
	\end{itemize}
	
	\noindent\textbf{Logros}
		\begin{itemize}
		\item Diseñé e implementé tableros de control en Elastic APM, con visualizaciones de KPIs de performance y negocio para cada uno de los procesos y servicios mantenidos por el equipo.
		\item % TODO: Alertas y SLOs
		\item Actué como referente técnico para validar iniciativas de alta complejidad, incluyendo el análisis de normativas impositivas internacionales y requerimientos de seguridad 
		\item % TODO: Diseño servicios y migración en 3 meses.
	\end{itemize}
	
	\cvExperience{Technical Lead - Geopagos}{Julio 2022 - Febrero 2025}{Ciudad Autónoma de Buenos Aires, Argentina}
	
	\noindent\textbf{Responsabilidades}
	
	\noindent\textbf{Logros}
	
	\cvExperience{Sr Backend Developer - Geopagos}{Octubre 2020 - Junio 2022}{Ciudad Autónoma de Buenos Aires, Argentina}
	
	\noindent\textbf{Responsabilidades}
	
	\noindent\textbf{Logros}
	
	\cvExperience{Fullstack Developer - Gloobes}{Abril 2014 - Septiembre 2020}{La Plata, Provincia de Buenos Aires, Argentina}
	
	\noindent\textbf{Responsabilidades}
	
	\noindent\textbf{Logros}





	\section{Formación académica} {
		\setlength{\parindent}{0pt}
		
		\textbf{Analista Programador Universitario:} Facultad de Informática, Universidad Nacional de La Plata. Graduado 2021.
	}
	
	
	\section{Idiomas} {
		\setlength{\parindent}{0pt}
		
		\textbf{Español:} Nativo.
		
		\textbf{Inglés:} Nivel técnico de lecto comprensión.
	}
	
\end{document}